\subsection{求解与结果}
基准时间窗下，我们先在 1.4--4.8 的粗网格和 3.0--3.4 的加密网格上扫描 $\lambda_{\max}$，得到首个连续正区间出现在 $K=3.2186$ 附近，并在 3.22--3.40 区间保持为正。表 \ref{tab:q4robust} 给出的局部复核表明，在总积分时长固定为 90 的前提下，$K=3.20$ 与 $K=3.22$ 之间的符号跨越对相邻步长和微扰初值仍可保留。

\begin{table}[htbp]
\centering
\caption{问题四短时间窗局部复核}
\label{tab:q4robust}
\small
\begin{tabularx}{0.95\textwidth}{>{\raggedright\arraybackslash}p{4.1cm} >{\centering\arraybackslash}p{1.9cm} >{\centering\arraybackslash}p{1.9cm} >{\centering\arraybackslash}p{1.9cm} X}
\toprule
设置 & $\lambda_{\max}(3.20)$ & $\lambda_{\max}(3.22)$ & 判别 & 说明 \\
\midrule
基准步长 $dt=0.010$ & $-0.00135$ & $0.00117$ & 符号跨越 & 3.22 附近进入正指数区间 \\
较小步长 $dt=0.008$ & $-0.00134$ & $0.00117$ & 符号跨越 & 结果与基准设置一致 \\
较大步长 $dt=0.012$ & $-0.00129$ & $0.00123$ & 符号跨越 & 结果与基准设置一致 \\
微扰初值 $(0.8005,0.1995,0.2)$ & $-0.00087$ & $0.00165$ & 符号跨越 & 微扰后仍保留同样的符号关系 \\
\bottomrule
\end{tabularx}
\normalsize
\end{table}

但是，我们一旦把积分时长继续拉长，正指数区间就会明显左移。总积分时长取 150 时，3.0--3.4 加密区间内各点均为正；总积分时长取 200 时，连续正区间从 3.11 开始，且 $\lambda_{\max}(3.20)=0.00579$、$\lambda_{\max}(3.22)=0.00701$，两点均已为正。进一步扩展到 5 组初值、5 个步长和 3 个积分时长的 75 组测试后，仅有 5 组（6.7\%）仍保留 $\lambda_{\max}(3.20)<0<\lambda_{\max}(3.22)$，即该符号跨越在扩展设置下不具备统计稳健性。

基准短时间窗下，$K=3.2186$ 附近确实出现了由负转正的局部跨越；把积分时长和数值设置继续放宽后，这个跨越又不能稳定保持，因此 $K=3.2186$ 不能作为唯一临界点。本文据此只保留“承载力提高会增强系统复杂性”这一方向性结论；若需精确定位分岔值，须以更长积分时长（$\ge 500$）配合分岔图交叉验证。图 \ref{fig:q4} 给出了不同积分时长下的有限时间指数扫描结果。

\begin{figure}[htbp]
\centering
\includegraphics[width=0.92\textwidth]{fig04_q4_lyapunov_scan.png}
\caption{问题四不同积分时长下的有限时间李雅普诺夫指数扫描}
\label{fig:q4}
\end{figure}

